% GENERATED FILE. Edit canonical lessons, then run npm run generate:books.
% canonical-source-hash: fnv1a64:0db96da661a3e819
% canonical-lessons: SA-C160-bhojanam, SA-C160-vyanjanam, SA-C160-aharah, SA-C160-asanam, SA-C160-kalasi

\chapter{Meal, Curry, Food, Eating, Small pitcher}
\label{ch:sa-meal-curry-food-eating-small-pitcher}
\hlchaptermodality{160}

\begin{chapteropening}
\textbf{By the end of this chapter:} \emph{I can say a meal, a curry, food, eating and a small pitcher.}

Complete the last lesson of chapter 160: I can say a meal, a curry, food, eating and a small pitcher.
\end{chapteropening}

\section[\texorpdfstring{\sk{भोजनम्} (bhojanam)}{भोजनम् (bhojanam)}]{\sk{भोजनम्} (bhojanam) — a meal, food}

\label{lesson:SA-C160-bhojanam}



\begin{quote}
Before the new one: say the Sanskrit for sesame, then the Sanskrit for a chickpea.
\end{quote}

\subsection*{You'll want to know: \sk{भोजनम्}}
\textbf{\sk{भोजनम्}} — \emph{bhojanam} — \textquotedblleft{}a meal, food\textquotedblright{}.

\begin{grammarlens}[title={What you've built}]
One more word for this chapter.
\end{grammarlens}

\subsection*{Your turn}
\begin{itemize}
\raggedright
  \item \emph{Say it:} \emph{bhojanam}
  \item \emph{Say it:} \emph{bhojanam}, once more
  \item \emph{Say it:} say \emph{bhojanam}
  \item \emph{Recall:} say the Sanskrit for sesame, then the Sanskrit for a chickpea, then say \emph{bhojanam} again
\end{itemize}

\begin{tcolorbox}[breakable,colback=teal!4,colframe=teal!35!black,title={Before you move on}]
What does \sk{भोजनम्} mean? (\textquotedblleft{}A meal, food\textquotedblright{}.) Say it once more.
\end{tcolorbox}
\section[\texorpdfstring{\sk{व्यञ्जनम्} (vyañjanam)}{व्यञ्जनम् (vyañjanam)}]{\sk{व्यञ्जनम्} (vyañjanam) — a curry, a side dish}

\label{lesson:SA-C160-vyanjanam}



\begin{quote}
Before the new one: say the Sanskrit for a chickpea, then the Sanskrit for a meal.
\end{quote}

\subsection*{You'll want to know: \sk{व्यञ्जनम्}}
\textbf{\sk{व्यञ्जनम्}} — \emph{vyañjanam} — \textquotedblleft{}a curry, a side dish\textquotedblright{}.

\begin{grammarlens}[title={What you've built}]
One more word for this chapter.
\end{grammarlens}

\subsection*{Your turn}
\begin{itemize}
\raggedright
  \item \emph{Say it:} \emph{vyañjanam}
  \item \emph{Say it:} \emph{vyañjanam}, once more
  \item \emph{Say it:} say \emph{vyañjanam}
  \item \emph{Recall:} say the Sanskrit for a chickpea, then the Sanskrit for a meal, then say \emph{vyañjanam} again
\end{itemize}

\begin{tcolorbox}[breakable,colback=teal!4,colframe=teal!35!black,title={Before you move on}]
What does \sk{व्यञ्जनम्} mean? (\textquotedblleft{}A curry, a side dish\textquotedblright{}.) Say it once more.
\end{tcolorbox}
\section[\texorpdfstring{\sk{आहारः} (āhāraḥ)}{आहारः (āhāraḥ)}]{\sk{आहारः} (āhāraḥ) — food, diet}

\label{lesson:SA-C160-aharah}



\begin{quote}
Before the new one: say the Sanskrit for a meal, then the Sanskrit for a curry.
\end{quote}

\subsection*{You'll want to know: \sk{आहारः}}
\textbf{\sk{आहारः}} — \emph{āhāraḥ} — \textquotedblleft{}food, diet\textquotedblright{}.

\begin{grammarlens}[title={What you've built}]
One more word for this chapter.
\end{grammarlens}

\subsection*{Your turn}
\begin{itemize}
\raggedright
  \item \emph{Say it:} \emph{āhāraḥ}
  \item \emph{Say it:} \emph{āhāraḥ}, once more
  \item \emph{Say it:} say \emph{āhāraḥ}
  \item \emph{Recall:} say the Sanskrit for a meal, then the Sanskrit for a curry, then say \emph{āhāraḥ} again
\end{itemize}

\begin{tcolorbox}[breakable,colback=teal!4,colframe=teal!35!black,title={Before you move on}]
What does \sk{आहारः} mean? (\textquotedblleft{}Food, diet\textquotedblright{}.) Say it once more.
\end{tcolorbox}
\section[\texorpdfstring{\sk{अशनम्} (aśanam)}{अशनम् (aśanam)}]{\sk{अशनम्} (aśanam) — eating, a meal}

\label{lesson:SA-C160-asanam}



\begin{quote}
Before the new one: say the Sanskrit for a curry, then the Sanskrit for food.
\end{quote}

\subsection*{You'll want to know: \sk{अशनम्}}
\textbf{\sk{अशनम्}} — \emph{aśanam} — \textquotedblleft{}eating, a meal\textquotedblright{}.

\begin{grammarlens}[title={What you've built}]
One more word for this chapter.
\end{grammarlens}

\subsection*{Your turn}
\begin{itemize}
\raggedright
  \item \emph{Say it:} \emph{aśanam}
  \item \emph{Say it:} \emph{aśanam}, once more
  \item \emph{Say it:} say \emph{aśanam}
  \item \emph{Recall:} say the Sanskrit for a curry, then the Sanskrit for food, then say \emph{aśanam} again
\end{itemize}

\begin{tcolorbox}[breakable,colback=teal!4,colframe=teal!35!black,title={Before you move on}]
What does \sk{अशनम्} mean? (\textquotedblleft{}Eating, a meal\textquotedblright{}.) Say it once more.
\end{tcolorbox}
\section[\texorpdfstring{\sk{कलशी} (kalaśī)}{कलशी (kalaśī)}]{\sk{कलशी} (kalaśī) — a small pitcher}

\label{lesson:SA-C160-kalasi}



\begin{quote}
Before the new one: say the Sanskrit for food, then the Sanskrit for eating.
\end{quote}

\subsection*{You'll want to know: \sk{कलशी}}
\textbf{\sk{कलशी}} — \emph{kalaśī} — \textquotedblleft{}a small pitcher\textquotedblright{}.

\begin{grammarlens}[title={What you've built}]
One more word for this chapter.
\end{grammarlens}

\subsection*{Your turn}
\begin{itemize}
\raggedright
  \item \emph{Say it:} \emph{kalaśī}
  \item \emph{Say it:} \emph{kalaśī}, once more
  \item \emph{Say it:} say \emph{kalaśī}
  \item \emph{Recall:} say the Sanskrit for food, then the Sanskrit for eating, then say \emph{kalaśī} again
\end{itemize}

\begin{tcolorbox}[breakable,colback=teal!4,colframe=teal!35!black,title={Before you move on}]
What does \sk{कलशी} mean? (\textquotedblleft{}A small pitcher\textquotedblright{}.) Say it once more.
\end{tcolorbox}
